\documentclass{article}
\usepackage[hypertexnames=false,colorlinks=true,breaklinks]{hyperref}
\usepackage{url}
\pagestyle{empty}		% no page numbers, thanks.
\addtolength{\topmargin}{-0.5in}% repairing LaTeX's huge margins...
% \setlength{\oddsidemargin}{0.5in}
\setlength{\oddsidemargin}{0.5in}
\addtolength{\textwidth}{1.5in}	% and here...
\addtolength{\marginparsep}{-0.5in}
\addtolength{\marginparwidth}{0.5in}
\setlength{\parindent}{0mm}	% indented paragraphs just don't make it
			   	% on a resume.  (My opinion; if you
			   	% disagree, give this whatever value you
			   	% want.)
\setlength{\textheight}{9.0in}	% more margin hacking.  I have a large
%\setlength{\textheight}{10.2in}	% more margin hacking.  I have a large
			    	% resume to fit on one page.

\begin{document}
\reversemarginpar		% this puts the margin notes on the
				% left-hand side of the resume.

{\huge{Joaqu\'{i}n Rapela}}	% make the name stand out

				% address stuff.  The \hspace*{\fill} is
				% necessary to get things laid out
				% correctly.  The sizes fed to \vspace are
				% variable, depending on how much
				% whitespace you want, and how much
				% stuff you're trying to fit on a page....
% \vspace{.1in}
\renewcommand{\labelitemi}{$\star$}
\textit{\textbf{Address}}: Gatsby Computational Neuroscience Unit. University
College London. 25 Howland St, Fitzrovia, London W1T 4JG, United Kingdom.
\textit{\textbf{Phone}}:~44 7432049398
\textit{\textbf{email}}:~j.rapela@ucl.ac.uk
\textit{\textbf{www}}:~http://www.gatsby.ucl.ac.uk/$\sim$rapela,
\textit{\textbf{github}}:~https://github.com/joacorapela,
\textit{\textbf{last update}}:~\today

%\par
% \vspace{2\baselineskip}
%\vspace{\baselineskip}
%A full-time position in software engineering or  \marginpar{{\bf Objective}}
%user interface design.
				% if the \marginpar is at the beginning
				% of the line, it loses.  *sigh*.  This
				% is a typical entry in a resume.   (I'm
				% defining an entry as something
				% important enough to have a marginal note
				% setting it off.)  Give it some space
				% with \vspace, use \par to break
				% between paragraphs, enter your text,
				% and place the \marginpar at the end.
\par
\vspace{2\baselineskip}
{\bf University of Southern California} \marginpar{{\bf Education}}

\par
\vspace{.05in}
{\bf Department of Electrical Engineering} \hspace*{\fill} 2003---2010
\par

\par
PhD in Electrical Engineering.
\par
Advisors: Dr. Norberto M. Grzywacz (Neuroscience) and Dr. Jerry M. Mendel (Signal Processing)
\par
\vspace{.05in}
{\bf Neuroscience Graduate Program} \hspace*{\fill} 2001---2003
\par

\par
Studies in Neuroscience.

\vspace{.05in}
{\bf Department of Electrical Engineering} \hspace*{\fill} 2000---2003

\par
MS in Electrical Engineering.

\par
\vspace{.1in}
{\bf Universidad de Buenos Aires}
\par
``Licenciatura (equivalent to MS)'' in Computer Science.
\hfill 1995---1998
\par
BS in Computer Science.\hfill 1992---1995

\par
\vspace{2\baselineskip}
{\bf Gatsby Computational Neuroscience Unit}
\marginpar{{\bf Academic}}
\par
{\bf University College London} \hspace*{\fill} June 2019---
\marginpar{{\bf Experience}}
\par
{\bf Position}: Research Engineer Fellow.

\par
\vspace{.1in}
{\bf The Truccolo Lab for Computational Neuroscience}
\par
{\bf Brown University} \hspace*{\fill} August 2017---January 2019
\par
{\bf Position}: Research Associate.

\par
\vspace{.1in}
{\bf Instituto en Luz Ambiente y Vision}
\par
{\bf Universidad Nacional de Tucum\'{a}n, Argentina} \hspace*{\fill} November 2013---Present
\par
{\bf Positions}: External researcher.

\par
\vspace{.1in}
{\bf Swartz Center for Computational Neuroscience}
\par
{\bf University of California San Diego} \hspace*{\fill} November 2010---July 2017
\par
{\bf Positions}: Postdoctoral and visiting scholar,

\par
\vspace{.1in}
{\bf University of Southern California}
\par
{\bf Position}: Research Assistant. \hspace*{\fill} 2001---2010
\par
{\bf Advisors}: Dr. Norberto M. Grzywacz. Director, Neuroscience Graduate
Program. Professor, Department of Biomedical Engineering. Dr. Jerry Mendel.
Professor, Department of Electrical Engineering.

\par
\vspace{.1in}
{\bf University of Southern California}
\par
{\bf Position}: Research Assistant. \hspace*{\fill} August 2000---August 2001
\par
{\bf Advisor}: Dr. Richard M. Leahy. Professor, Department of Electrical
Engineering.

\par
\vspace{2\baselineskip}
{\bf Brown University} \hspace*{\fill} September 2017---December 2017
\marginpar{{\bf Teaching}}

\par
{\bf Course}: NEUR2110 Statistical Neuroscience
\marginpar{{\bf Experience}}

\par
{\bf Role}: Graded homeworks and addressed questions from students related to course material.

\par
\vspace{.1in}
{\bf University of Southern California} \hspace*{\fill} January 2009---May 2009
\par
{\bf Course}: EE 364: Introduction to Probability and Statistics for Electrical Engineering and Computer Science. Evaluations available upon request.
\par
{\bf Role}: Discussion section leader.

\par
\vspace{.1in}
{\bf Universidad de Buenos Aires} \hspace*{\fill} August 1999---December 1999
\par
{\bf Course}: Numerical Linear Algebra.
\par
{\bf Role}: Teaching Assistant.

\par
\vspace{.1in}
{\bf Universidad de Buenos Aires} \hspace*{\fill} August 1998---August 1999
\par
{\bf Course}: Artificial Intelligence.
\par
{\bf Role}: Teaching Assistant.

\par
\vspace{2\baselineskip}
{\bf Vision Research, Frontiers in Neuroscience, Journal of Perceptual Imaging}
\marginpar{{\bf Service}}
\par
Ad hoc reviewer.
\par

\par
\vspace{2\baselineskip}
{\bf Center for Brain Activity Mapping}  \hspace*{\fill} 2013---2014
\marginpar{{\bf Funding }}
\par
Principal investigator, grant \texttt{2013-023CBAM}
\par
\emph{Taking the next step toward understanding computations by neural
ensembles with high resolution neural recordings, generic data assimilation
methods, and increased computational power}.

\par
\vspace{2\baselineskip}
{\bf University of Southern California}
\marginpar{{\bf Selected}}
\par
\marginpar{{\bf Courses}}
{\bf Mathematics:} Topology, Fundamental Concepts of Analysis, Real Analysis (audited), Functional Analysis (audited). {\bf Engineering:} Transform Theory for Engineers, Random Processes in Engineering, Introduction to Digital Signal Processing, Computational Solution to Optimization Problems, Information Theory, Statistics for Engineers. {\bf Neuroscience:} Control and Communication in the Nervous System, Neurobiology, Advanced Neurosciences I, Advanced Neurosciences II.\par
\par

\vspace{2\baselineskip}
{\bf IBM Almaden Research Center}
\marginpar{{\bf Industry}}
\par
{\bf Position}:  Staff Software Engineer
\hspace*{\fill} January 2000---August 2000
\marginpar{{\bf Experience}}

\par
\vspace{.1in}
{\bf Oracle, Argentina}
\par
{\bf Position}:  Software Engineer  \hfill June 1998---September 1998

\par
\vspace{.1in}
{\bf IBM, Argentina}
\par
{\bf Position}:  Fellowship holder  \hfill September 1997---June 1998
\par

\vspace{.1in}
{\bf IBM Almaden Research Center}
\nopagebreak
\par
{\bf Position}:  Fellowship holder  \hfill October 1996---September 1997
\nopagebreak
\par

\par
\vspace{.1in}
{\bf IBM, Argentina}
\par
{\bf Position}:  Fellowship holder  \hfill April 1994---October 1996
\par

\vspace{2\baselineskip}
R, Python, Matlab, Java, C/C++, Smalltalk, Pascal, Assembler.
\marginpar{{\bf Programming}}
\par
\marginpar{{\bf Languages}}
\par

\end{document}

